% Einheit 1 von 5 – Rechteck, Quadrat, Umfang (bank/flaechen/e1.jsonl, zusammenbau v0.1)
%% TODO Zweigzeile Teil 1: was hier gelernt wird (Satz in Schülersprache aus dem Einheitstitel) – Einheit 1
\einheitenkopf[e1]{Einheit 1 von 5 · Rechteck, Quadrat, Umfang}
\zweigzeile{ab Kl. 5 · P10 oft}
\setcounter{aufgabe}{9}

%% TODO Titel: Ich-kann-Satz fehlt in der Bank (Platzhalter: Kettenname) – Nr. 10
%% TODO Anweisung: ein Satz über der ersten Teilaufgabe fehlt in der Bank – Nr. 10
\begin{aufgabe}{Gegeben – gesucht}
\begin{teile}
\teil Ein Rechteck hat die Fläche $A = 32$ m² und die Breite $b = 4$ m. Was ist gesucht? Kreise es in $A = a \cdot b$ ein.
\end{teile}
\end{aufgabe}

%% TODO Titel: Ich-kann-Satz fehlt in der Bank (Platzhalter: Kettenname) – Nr. 11
%% TODO Anweisung: ein Satz über der ersten Teilaufgabe fehlt in der Bank – Nr. 11
\begin{aufgabe}{Rechteck}
%% TODO \rechenplatz{n}: Schrittzahl des Grundfalls fehlt in der Bank (nur bei fester Schreibform, 2.3 b)
\begin{teile}
\teil Ein Garten bekommt einen Zaun. Fläche oder Umfang? \\ \kreuz{Fläche} \\ \kreuz{Umfang}
\teil Rechteck, $a = 8$ cm, $b = 3$ cm – Fläche und Umfang? A = \leerfeld[cm²] , u = \leerfeld[cm]
\teil Rechteck, $a = 4{,}5$ cm, $b = 2$ cm – Fläche und Umfang? A = \leerfeld[cm²] , u = \leerfeld[cm]
\teil Quadrat, $a = 6$ cm – Fläche und Umfang? A = \leerfeld[cm²] , u = \leerfeld[cm]
\teil Ein Fünfeck hat die Seiten $3$ cm, $4$ cm, $5$ cm, $2$ cm und $6$ cm – Umfang? u = \leerfeld[cm]
\teil Rechteck, $A = 56$ cm², $a = 8$ cm – wie lang ist $b$? b = \leerfeld[cm]
\teil Rechteck, $u = 20$ cm, $a = 6$ cm – wie lang ist $b$? b = \leerfeld[cm]
\teil Quadrat, $A = 64$ cm² – Seitenlänge? a = \leerfeld[cm]
\teil Ein Rechteck ist doppelt so lang wie breit; die Breite ist $x$. Term für den Umfang? u = \leerfeld
\teil Ein rechteckiger Bilderrahmen hat einen Umfang von $34$ cm, eine Seite ist $11$ cm lang. Wie lang ist die andere Seite? \hfill (P10 2018 OS) \\ \leerfeld[cm]
\end{teile}
\end{aufgabe}

%% TODO Titel: Ich-kann-Satz fehlt in der Bank (Platzhalter: Kettenname) – Nr. 12
%% TODO Anweisung: ein Satz über der ersten Teilaufgabe fehlt in der Bank – Nr. 12
\begin{aufgabe}{aus Rechtecken zusammengesetzte Fläche (Summe)}
\begin{teile}
\teil Ein L-förmiger Raum besteht aus zwei Rechtecken: $6$ m × $4$ m und $3$ m × $2$ m. Fläche? \leerfeld[m²]
\end{teile}
\end{aufgabe}

%% TODO Titel: Ich-kann-Satz fehlt in der Bank (Platzhalter: Kettenname) – Nr. 13
%% TODO Anweisung: ein Satz über der ersten Teilaufgabe fehlt in der Bank – Nr. 13
\begin{aufgabe}{Einheit wechseln vor dem Rechnen (cm und m gemischt)}
\begin{teile}
\teil Rechteck, $2$ m lang, $45$ cm breit – Fläche in m²? \leerfeld[m²]
\end{teile}
\end{aufgabe}

%% TODO Titel: Ich-kann-Satz fehlt in der Bank (Platzhalter: Kettenname) – Nr. 14
%% TODO Anweisung: ein Satz über der ersten Teilaufgabe fehlt in der Bank – Nr. 14
\begin{aufgabe}{Rechteck – Fehler finden, Begründen, Anwendung, Darstellungswechsel}
\begin{teile}
\teil Ein Rechteck ist $6$ cm lang und $2{,}5$ cm breit. Tom rechnet den Umfang: \rechnung{u &= 6 \cdot 2{,}5 \\ u &= 15 \text{ cm}} Wo ist der Fehler?
\teil Warum rechnet man beim Rechteck Länge mal Breite? Erkläre mit Einheitsquadraten.
\teil Ein Zimmer ist $4{,}2$ m lang und $3{,}5$ m breit. Laminat gibt es in Paketen zu $2{,}2$ m². Wie viele Pakete muss man kaufen?
\teil Zeichne ein Rechteck, dessen Fläche $A = 3 \cdot x$ ist, und beschrifte die Seiten.

\rechenplatz[halb]{4}
\end{teile}
\end{aufgabe}
